%!TEX root = rapport.tex

\chapter{Representation of Lambda Terms in Python}

A lambda term is represented as a directed acyclic graph (DAG). The nodes of the
DAG are encoded in Python as objects. There are three types of nodes: 
abstraction, application and variable:
\begin{description}
	\item[Abstraction nodes] define scopes for bound variables and can be thought
	of as the encoding of the $\lambda$s in the DAG-structure. This node type
	always has one child.
	
	\item[Application nodes] have two children. The first of these
	is either the DAG-representation of some \emph{active} subterm of 
	$M$, or another application node. The second child is the DAG-representation
	of the ``argument'', that is, the lambda term that are to be substituted
	in place for a bound variable. An application node represents a possible
	one-step $\beta$-reduction.
	
	\item[Variable nodes] represent free or bound variables. A bound 
	variable is associated with a scope, represented by an ancestral
	abstraction node. Variable nodes cannot have descendants.
\end{description}

This representation is straightforwardly encoded in Python with a general node
class, \pycode{LambdaNode}, which has a subclass for each node type.
\begin{itemize}
	\item Each \pycode{LambdaNode}-object maintains pointers to its immediate 
	descendants as well as a pointer to its parent node.
	
	\item The abstraction nodes furthermore maintain a list of pointers to the
	variable nodes that represent the variables bound by this particular abstractor.
	These variable nodes must be descendants of the abstraction node, although 
	not necessarily direct descendants.
	
	\item Every abstraction node knows the name of the variable that it binds.
	This is done in order to be able to easily build the list of bound 
	variables after parsing a string representation of a lambda term. When
	the list of variables has been built, this name is not important anymore, 
	but it still serves as a means of producing human-readable output of the DAG.
	
	\item When a term is reduced, the problem is attacked ``top-down'', i.e. we first
	see the abstraction node, and then we have to deal with its variables.
	Therefore there is no need for a doubly linked connection between abstraction
	nodes and variable nodes, and we can safely encode the relationship as a simple
	list of bound nodes in each abstraction node.
\end{itemize}

A term that is not in $\beta$-normal-form has at least one $\beta$-contractum,
and therefore the DAG-representation of that term contains at least one
application node and one abstraction node. More specifically, for each
possible one-step $\beta$-reduction in a term $M$, its DAG will contain one
application node. The first descendant of this node will be an abstraction
node, while the second, ``rightmost'', descendant will be the argument that
will be applied to the function represented by the abstraction node.

Given this DAG, it is possible to implement the aforementioned operations on lambda
terms by treewalks. 


% Callback functions? Hvad skal de bruges til?


% \x.xy bliver til
%    L
%    |
%    @
%   / \ 
%  x   y


% application node: @
% abstraction node: L
% variable node: x,y..

% identify each of the possible reducts: find all active subterms.


% (\x.xx)(\y.\z.zyz) == (\x.xx)(\y.\z.(zy)z)
%        @
%       / \
%      L   ------L
%      |         |
%      @         @
%     / \       / \
%    x  x      @   z
%             / \
%            z   y
% .... or:
%       @
%      / \
%    Lx   Ly
%     |   |
%    xx   Lz
%         |
%         zyz

% (\x.xx)(\y.yy)(\z.zz) == ((\x.xx)(\y.yy)) (\z.zz)
%          @
%         / \
%        @   Lz
%       / \  |
%     Lx  Ly zz
%      |   | 
%     xx  yy 
% .... becomes:
%       @
%      / \
%     @   --- Lz
%    / \      |
%   Ly  Ly    zz
%   |   |
%   yy  yy
%

% I(Iz) -> Iz
% (\x.x)((\y.y)z) -> (\y.y)z
%          @
%         / \
%        Lx  ---@
%        |     / \
%        x    Ly  z
%             |
%             y
% indeholder to one-step betareductions:
%     @
%    / \
%   Ly  z
%   |
%   y
% og (ækvivalent!):
%     @
%    / \
%   Lx  z
%   |
%   x
